\documentclass[instructions.tex]{subfiles}
\begin{document}
 
\chapter{Du respect que l’homme doit avoir pour son corps.}
 
\primo Un prince qui place un sujet dans un palais enrichi de tout ce qu’il y a de plus magnifique et de plus délicieux, fait connoître combien ce sujet lui est cher. O homme ! comprenez combien vous êtes cher à Dieu, combien il vous estime ; c’est pour vous qu’il a créé l’univers, mille fois plus magnifique que les Louvres des rois. Qu’est-ce donc que l’homme, ô mon Dieu ! disoit le Prophète, pour que vous daigniez penser à lui ? Vous l’avez presque égalé aux anges ; vous l’avez couronné comme un roi dans l’univers : tout ce qu’il y a dans la mer, dans les airs et sur la terre, vous lui avez tout soumis\footnote{Ps. 8.}.
 
Quel respect ne devons-nous pas avoir pour nous-mêmes, puisque Dieu nous honore de la sorte ! Et avec quel respect devons-nous traiter les autres hommes ! Les rois n’ont-ils pas du respect pour les autres souverains ? Ne traitent-ils pas avec honneur les princes subalternes ? Tous les hommes étant en un sens comme rois dans l’univers, doivent donc se prévenir par des marques d’honneur et d’estime, selon l’ordre et la subordination des différentes conditions\footnote{Honore invicem prævenientes. Rom. 12. 10.}.
 
\secundo Si l’homme est si respectable par rapport aux objets qui sont hors de lui, combien l’est-il plus en lui-même, soit par rapport à son corps, soit par rapport à son âme ! Son corps est, de toutes les créatures visibles, l’ouvrage le plus beau et le plus parfait qui soit sorti des mains de Dieu. Si le corps souillé par le péché mérite du mépris et des châtimens, il mérite aussi nos respects, si on le considère dans sa création et dans les desseins de Dieu.
 
Le Créateur, ayant de toute éternité résolu d’envoyer son Fils sur la terre ; et de lui donner un corps capable des plus nobles opérations, a formé notre corps à l’image du corps adorable de l’Homme-Dieu, qui est comme notre frère aîné, notre prototype et notre original. Voilà quelle est la dignité de notre origine, selon le corps. En comprenez-vous la noblesse ? Vous devez le traiter avec respect et avec honneur ; pourquoi l’avilissez-vous par une conduite indigne de ce que vous êtes ?
 
Glorifiez Dieu dans votre corps, dit saint Paul\footnote{Glorificate et portale Deum in corpore vestro. 1. Cor. 6.20.} : c’est-à-dire, que tous vos sens soient employés à des actions dignes de votre auteur; vos yeux à contempler ses ouvrages, votre langue à publier ses grandeurs, vos oreilles à entendre ses louanges, vos facultés à agir pour sa gloire, et vivre selon ses desseins pour le service de vos frères; c’est ainsi que vous ferez de vos corps des victimes saintes. Mais, hélas ! qu’en faites-vous ?
 
\tertio L’homme, si digne de respect dans l’ordre naturel, l’est incomparablement plus dans l’ordre surnaturel. Destinés à régner dans la vie future, tous les hommes qui l’auront mérité, seront autant de rois couronnés de gloire. Leurs corps ressuscités seront doués des plus glorieuses qualités, d’un éclat, d’une agilité, d’une subtilité surprenantes; ils seront impassibles et immortels. Oserions-nous nous plaindre de quelques peines que nous devons prendre pour procurer à notre corps un bonheur si grand ?
 
O mon corps, que vous êtes noble, que la fin à laquelle vous êtes destiné est auguste ! vous êtes l’ouvrage de Dieu : vous lui appartenez plus qu’à moi. C’est par le respect que j’ai pour vous, c’est pour ne pas vous déshonorer, et pour vous rendre heureux, que je dois vous exercer par les travaux d’une vie chaste et sainte. Si vous souffrez ici-bas, vous serez un jour glorifié.
 
\section{Résolutions}
 
\primo Je respecterai, dans mon prochain, et en ma personne, l’image de Dieu.
 
\secundo J’honorerai dans mon corps la représentation de celui de Jésus-Christ. Et je le sanctifierai par la mortification de ses convoitises, par les exercices du travail et de la pénitence. 

\prayer{Faites, ô mon Dieu ! que je gouverne désormais si sagement mon corps, qu’il ressuscite glorieusement à la fin du monde.}
 
\end{document}
